\documentclass{article}
\usepackage{graphicx} % Required for inserting images
\usepackage{caption} % Required for \captionof
\usepackage{amsmath}
\usepackage{float}
\usepackage{comment}
\usepackage[T1]{fontenc}
\usepackage[utf8]{inputenc}
\usepackage{lmodern}
\usepackage{microtype}
\usepackage{xcolor}
\usepackage{framed}
\renewenvironment{leftbar}{%
  \def\FrameCommand{\textcolor{orange}{\vrule width 3pt} \hspace{10pt}}%
  \MakeFramed {\advance\hsize-\width \FrameRestore}}%
 {\endMakeFramed}

\title{Mechanikai mérés}
\author{Kern Luca, Vincze Csongor}
\date{2026 Április 22}

\begin{document}

\maketitle

\section*{8. Feladat: Kényszerrezgés amplitúdójának és sebességamplitúdójának vizsgálata a kényszerítő frekvencia függvényében.}

Kis csillapítás: 
Csillaító mágneseket visszaaszereltük, a távolságuk:
$d_{\text{mágnesek}} = 2,0 \text{cm}$
$m_{\text{teljes}} = 100 \text{g}$
Impulzussorozat beállításai (1-es csatorna): 6400 Hz, 3 $\mu\text{s}$ , 3 $V_{\text{pp}}$
2-es csatorna: 3 V, DC
Kijelzőn: 5-10 impulzus látható

\begin{table}[H]
    \centering
    \begin{tabular}{|c|c|}
        \hline
        Frekvencia [Hz] & Amplitúdó [mm] \\ \hline
          0 & \\ \hline 
          25 & \\ \hline
          50 & \\ \hline 
          75 & \\ \hline
          100 & \\ \hline 
          125 & \\ \hline
        
    \end{tabular}
    \caption{Amplitúdó a frekvencia függvényében}
    \label{tab:amplitudo}
\end{table}

Rezonanciafrekvencia: $f_{\text{rez}} = \text{Hz}$

Nagy csillapítás:

A mágnesek távolságát csökkentettük: $d_{\text{mágnesek}} = \text{cm}$
Impulzussorozat beállításai (1-es csatorna): 6400 Hz, 3 $\mu\text{s}$ , 3 $V_{\text{pp}}$
2-es csatorna: 3 V, DC
Kijelzőn: 5-10 impulzus látható

\begin{table}[H]
    \centering
    \begin{tabular}{|c|c|}
        \hline
        Frekvencia [Hz] & Amplitúdó [mm] \\ \hline
          0 & \\ \hline 
          25 & \\ \hline
          50 & \\ \hline 
          75 & \\ \hline
          100 & \\ \hline 
          125 & \\ \hline
        
    \end{tabular}
    \caption{Amplitúdó a frekvencia függvényében}
    \label{tab:amplitudo}
\end{table}

Rezonanciafrekvencia: $f_{\text{rez}} = \text{Hz}$

\end{document}